\scanpage{529}%
\mspage{440}%
% archive sheet number: 528
de façon \unsure{ess\textsuperscript{t}.|esst.} unique par un $\lambda_i^*(y_i)$.

Cette fois par contre, c'est le cas cofibrant qui pose problème. En fait, je n'arrive pas à m'en tirer.

Ainsi, on trouve bien un foncteur

\[ \mathrm{Ind}_\pi\bigl(\unsure{E|\mathcal E}(\pi)\bigr)\longrightarrow E, \]

qui est ess. surjectif, mais faute de savoir que $E(\pi)\subset E_\pi$ dans le cas cofibrant, il n'est pas clair que ce foncteur soit une équivalence de catégories.

4.15.7. Il me faut \unsure{reporter|reprendre} la démonstration du lemme 4.15.3\inserted{(p. \unsure{428|420})}, en admettant 4.8.13. (\unsure{Non encore démontré|Non encore démontrée} -- Je vais revenir plus bas sur la dém. de 4.8.13), On va utiliser le

\uline{Lemme} 4.15.7.1. Soi\deleted{\unsure{ent|nt}}t $C$ une catégorie, $\underline{x}=(x_i)_{i\in I}$ un Ind-objet de $C$. \deleted{\unsure{a)|g}}\deleted{qui soit}\inserted{Pour} que \deleted{$\unsure{C^\natural|C}$ une enveloppe de Karoubi de $C$, $x_0$ un objet de $C^{\deleted{\unsure{\natural|E}}}$ tel que \deleted{\unsure{qui|que}} $\underline{x}$ soit isomorphe dans $\mathrm{Ind}(C)$,} $\underline{x}$ soit \inserted{ess. constant, i.e.} isomorphe, dans $\mathrm{Ind}(C)$, à un \deleted{\unsure{un|m}} Ind-objet constant \deleted{\unsure{et|d}} défini par un objet \deleted{\illegible{1}}
